
\section{Related Work}


RAG~\citep{guu2020retrieval,lewis2020retrieval,khattab2021relevance,izacard2022few}) is now a common strategy for bolstering LLMs by combining them with retrieval systems. Through retrieval, RAG helps LM systems gather domain-specific knowledge, ground generations in factual information~\cite{shuster2021retrieval, huo2023retrieving}, and offer a degree of transparency or interpretability via citing sources~\cite{mialon2023augmented}.



Multiple LLM-based evaluation techniques have emerged for gauging LLM systems. This is essential for rapid deployment in new settings, where it is difficult to build a traditional benchmark dataset from scratch.
Early attempts at this use LLMs out of the box, as in MT-Bench and Chatbot Arena~\cite{zheng2023judging}.
AutoCalibrate~\cite{liu2023calibrating} seeks to align an LLM-judge with human preferences, leveraging a self-refinement prompt to iteratively improve the LLM judge. 
However, AutoCalibrate does not offer any statistical guarantees for the accuracy of its predictions.
Other work has used LLM prompting to evaluate system quality across natural language generation tasks, such as translation, summarization, and dialogue  \cite{kocmi2023large,fu2023gptscore, liu2303g, wang2023chatgpt}.

In the context of knowledge-intensive NLP tasks, LLMs have been explored for assessing attribution and factuality in LLMs \cite{min2023factscore, gekhman2023trueteacher, yue2023automatic}.
New guidelines like LongEval \cite{krishna-etal-2023-longeval} and datasets like Hagrid and ALCE \cite{kamalloo2023hagrid, gao2023enabling} provide resources for analyzing knowledge-intensive LLM pipelines.


The two most-closely related projects to ARES are EXAM~\cite{sander2021exam} and RAGAS~\cite{ragas2023}. To evaluate RAG systems, the EXAM metric estimates how many exam questions a reader (simulated as a QA system) can answer correctly based on the generated response. This requires a set of queries with several associated sub-questions each, which adds a burden that ARES does not bring. RAGAS is based on a handful of heuristic hand-written prompts. These offer little adaptability to new RAG evaluation settings (e.g., new corpora) and, as we show in our evaluation, substantially underperform ARES.









